\documentclass[11pt]{article}
\usepackage[a4paper,margin=27mm]{geometry}
\usepackage{amsmath,amssymb,amsthm,booktabs,longtable,array}
\usepackage[T1]{fontenc}
\usepackage{lmodern}
\usepackage[colorlinks=true,linkcolor=blue,urlcolor=blue,citecolor=blue]{hyperref}
\hypersetup{pdftitle={Alternative proofs of Hohn's divisibility conjectures after Kontorovich and Stong},pdfauthor={Haobo Ma; Wenlin Zhang}}
\newtheorem{theorem}{Theorem}
\newtheorem{lemma}[theorem]{Lemma}
\newtheorem{corollary}[theorem]{Corollary}
\newcommand{\vp}{v_p}
\title{Alternative proofs of Hohn's divisibility conjectures\\
after Kontorovich and Stong}
\author{Haobo Ma$^1$ \and Wenlin Zhang$^{2,1}$}
\date{18 November 2026}
\begin{document}
\maketitle
\begin{center}
\small $^1$ChronoAI Pte Ltd, Singapore\\
$^2$National University of Singapore, Singapore
\end{center}
\begin{abstract}
For Charles L. Hohn's conjectures in OEIS A390148, revision 23
(17 November 2025), COMMENTS, paragraphs 5--9, we give alternative
proofs of the divisibility assertions using arithmetic already present in
Kontorovich (2012, Lemma 2.11) and Stong (2015, p.~76)
\cite{kontorovich,stong}, and prove the repeated-radius count.
Every positive prime power dividing one radius divides exactly one or two
other coordinates. The conclusion holds for every
prime, not just primes congruent to one modulo three. At the prime three the
radius valuations are $(0,e,e,e)$ with $e>0$; at any prime congruent to two
modulo three that occurs, they are $(0,2m,2m,2m)$ with $m>0$, up to
permutation. The same argument proves the prime-support assertion in Hohn's
gcd formula for the fourth radius. For the linked circle
array, an elementary parametrization proves the cube-or-twice-a-cube
difference conjecture. All these statements have complete written proofs.
\end{abstract}

\section{The equation and the result}
Hohn's array \cite{hohn} consists of ordered quadruples of positive integer
radii with common gcd one satisfying
\begin{equation}\label{eq:radii}
 \left(\sum_{i=1}^4\frac1{r_i}\right)^2
       =3\sum_{i=1}^4\frac1{r_i^2}.
\end{equation}
Here ``common gcd one'' does not mean pairwise coprime. Equation
\eqref{eq:radii} is the Soddy--Gossett equation in three dimensions, with
the curvature of a tangent plane set to zero; see \cite[Theorem 1.2]{lmw}.
Our arguments concern all positive integer solutions of this equation,
without requiring an ordering or distinct coordinates.

Write $v_p(n)$ for the exponent of the prime $p$ in a positive integer $n$.
When differences of integers occur, set $v_p(0)=+\infty$ and
$v_p(n)=v_p(|n|)$ for $n\ne0$. Each coordinate is counted separately,
including coordinates with equal radii.

\begin{theorem}\label{thm:main}
Let $r_1,r_2,r_3,r_4$ be positive integers satisfying \eqref{eq:radii}
and having common gcd one. For any prime $p$, put
$E=\max_i v_p(r_i)$.
\begin{enumerate}
\item If $p\equiv1\pmod3$ and $E>0$, the maximum $E$ occurs at least
twice, and $\min_i v_p(r_i)=0$.
\item For $p=3$, one has $E>0$ and the valuations are $(0,E,E,E)$,
up to permutation.
\item If $p\equiv2\pmod3$ and $E>0$, including $p=2$, the valuations
are $(0,E,E,E)$ up to permutation, and $E$ is even.
\end{enumerate}
Consequently, for every prime $p$, integer $k\ge1$, and coordinate $i$
with $p^k\mid r_i$, exactly one or two indices $j\ne i$ satisfy
$p^k\mid r_j$.
\end{theorem}

The last sentence is Hohn's split-prime assertion with its prime restriction
removed. The additional valuation conclusions address the adjacent comments
for $3$ and for primes congruent to $2$ modulo $3$. They concern every
positive exponent, rather than a finite initial segment of the array.
The split-prime clause says ``1 or 2 of the other terms'' at
\cite[COMMENTS, paragraph 6]{hohn}; paragraphs 5 and 7 give the other
two valuation assertions. The corresponding lines in the
\href{https://raw.githubusercontent.com/oeis/oeisdata/d4ade8b4ca7c6ff0ac8931532918fd4f6176805e/seq/A390/A390148.seq}{fixed sequence text}
are 10--12. Paragraph 8 (line 13) gives the repeated-radius count below.

\paragraph{Earlier arithmetic.}
Kontorovich's primitive Soddy-quintuple classification modulo $3$
\cite[Lemma 2.11, p.~9]{kontorovich}, first released in 2012, applies
to $(0,b_1,b_2,b_3,b_4)$ in \eqref{eq:curvatures}. It implies the
three-unit curvature pattern and hence our $3$-adic radius assertion.
McGonagle and Northshield's earlier article \cite[Section 6]{mcgonagle}
states the norm conclusion in a three-dimensional sketch.
Stong's January 2015 solution of Northshield's problem
\cite[p.~76]{stong} supplies the inert-prime norm argument;
Northshield gave a geometric proof in March 2015
\cite[Corollary 4.6, p.~14]{northshield}.
These are earlier results in equivalent arithmetic variables, not
assertions about Hohn's later sequence. We give the translation explicitly
and retain a direct proof. The sphere and circle repetition classifications
and the finite membership consequence are additional results here;
no priority is claimed.

\section{Integral curvatures and the proof}
Put
\begin{equation}\label{eq:curvatures}
 L=\mathop{\rm lcm}(r_1,r_2,r_3,r_4),\qquad b_i=L/r_i.
\end{equation}
These are positive integers, and clearing denominators in
\eqref{eq:radii} gives
\begin{equation}\label{eq:integral}
 \left(\sum_{i=1}^4 b_i\right)^2=3\sum_{i=1}^4 b_i^2.
\end{equation}
Their common gcd is one. Indeed, if $d$ divides every $b_i$, then
$L/d=r_i(b_i/d)$ is an integer common multiple of all the $r_i$.
Minimality of $L$ forces $d=1$. For every prime $p$,
\begin{equation}\label{eq:orders}
 v_p(b_i)+v_p(r_i)=v_p(L)=E.
\end{equation}
In particular, a curvature not divisible by $p$ corresponds to a radius of
maximal $p$-adic order. Primitivity of the radii gives
$\min_i v_p(r_i)=0$.

\begin{lemma}\label{lem:norm}
If $p\equiv2\pmod3$ is prime, then
$x^2-xy+y^2=0$ in $\mathbb F_p$ implies $x=y=0$.
For integers $x,y$ not both zero, it follows that
\[
 v_p(x^2-xy+y^2)=2\min\{v_p(x),v_p(y)\}.
\]
\end{lemma}
\begin{proof}
If $y=0$, the equation gives $x=0$. Otherwise let $z=-x/y$ in
$\mathbb F_p$. Then $z^2+z+1=0$, so $z^3=1$, whereas $z\ne1$
because $p\ne3$. This would give an element of order $3$ in the group
$\mathbb F_p^\times$ of order $p-1$, impossible when $p\equiv2\pmod3$.
This also covers $p=2$. For the integer assertion, factor out
$p^{2s}$, where $s=\min\{v_p(x),v_p(y)\}$. The remaining two arguments
are not both zero modulo $p$, so their norm is not divisible by $p$.
\end{proof}

\begin{proof}[Proof of Theorem \ref{thm:main}]
For an odd prime $p$, the primitive curvature quadruple cannot have exactly
one coordinate nonzero modulo $p$. If that coordinate is $b_i$,
\eqref{eq:integral} would give $b_i^2=3b_i^2$, or $2b_i^2=0$,
in $\mathbb F_p$, a contradiction. Thus at least two curvatures are
units modulo $p$, proving the first assertion by \eqref{eq:orders}.

At $p=3$, equation \eqref{eq:integral} first implies
$3\mid\sum b_i$. Writing $\sum b_i=3h$ then gives
$\sum b_i^2=3h^2$. A nonzero square modulo $3$ equals $1$.
Hence the number of curvatures not divisible by $3$ is a positive
multiple of $3$ not exceeding $4$, and is therefore $3$.
By \eqref{eq:orders}, exactly three radii have valuation $E$.
The fourth radius has valuation zero, since otherwise the radii would
have common factor $3$. Also $E>0$: if $E=0$, all four curvatures
would be units modulo $3$.

Now suppose $p\equiv2\pmod3$ and $E>0$. Relabel so that
$t=b_1$ has valuation $E$, which is possible by primitivity of the
radii and \eqref{eq:orders}. Some other curvature, say $a=b_2$,
is a unit modulo $p$. Write $b_3=a+u$, $b_4=a+v$, and
$S=b_2+b_3+b_4=3a+u+v$. Direct expansion of
\eqref{eq:integral} yields
\begin{equation}\label{eq:norm}
 u^2-uv+v^2=t(S-t).
\end{equation}
Reducing modulo $p$ and applying Lemma \ref{lem:norm} shows
$u\equiv v\equiv0\pmod p$. Thus the other three curvatures are units,
and $S-t\equiv3a\not\equiv0\pmod p$. The right side of
\eqref{eq:norm} has valuation exactly $E$. It is nonzero, and
Lemma \ref{lem:norm} makes its valuation even. Formula
\eqref{eq:orders} now gives the third assertion, including at $p=2$.

Finally, whenever $p^k$ divides a radius, $1\le k\le E$.
The maximum valuation occurs at least twice in each of the three cases,
so at least two coordinates meet this threshold. Primitivity excludes
all four coordinates meeting a positive threshold. The number meeting it
is therefore $2$ or $3$, leaving exactly one or two indices other than $i$.
\end{proof}

\paragraph{Translation of the earlier norm theorem.}
Write $\sum b_i=3h$ and set $c_i=h-b_i$. Then
$\sum c_i=h$, $\sum c_i^2=h^2$, and
$b_i=\sum_{j\ne i}c_j$. A common divisor of the $c_i$ divides $h$
and every $b_i$, so their common gcd is one. Stong's theorem, also
Northshield's Corollary 4.6, makes each positive $b_i$ an Eisenstein
norm, whose valuation at $p\equiv2\pmod3$ is even by Lemma
\ref{lem:norm}. Moreover Stong's displayed identity reads
$b_3b_4=c_1^2+c_1c_2+c_2^2$. If $p\mid b_3$, anisotropy forces
$c_1=c_2=0$ modulo $p$; the relation
$b_3=c_1+c_2+c_4$ then gives $c_4=0$ modulo $p$.
Primitivity forces $c_3\ne0$, so $b_1,b_2,b_4$ are units.
This proves exactly the inert-prime pattern via the earlier theorem,
including at $2$, without attributing the radius formulation to it.

Both threshold counts occur even at a split prime. For $p=7$, the
quadruple $(3,3,7,21)$ has two divisible coordinates, and
$(12,21,28,84)$ has three. The quadruple $(1,3,3,3)$ illustrates why
equal radii must be counted as separate coordinates.

\section{Consequences for the same array}
\begin{corollary}\label{cor:gcd}
For any labeling of a quadruple in Theorem \ref{thm:main}, the number
\[
 x=\frac{r_1r_2r_3}
 {r_4\gcd(r_1,r_2)\gcd(r_1,r_3)\gcd(r_2,r_3)}
\]
is a positive integer. Every prime factor of $x$ is congruent to $1$
modulo $3$ and divides at least two of $r_1,r_2,r_3$.
\end{corollary}
\begin{proof}
Fix a prime $p$, and sort the valuations of the first three radii as
$\alpha\le\beta\le\gamma$. Put $\delta=v_p(r_4)$. The exponent of
$p$ in the displayed rational number is
\[
 (\alpha+\beta+\gamma)
 -(\alpha+\alpha+\beta)-\delta=\gamma-\alpha-\delta.
\]
If $\delta=0$, this is nonnegative. If $\delta>0$, primitivity
forces $\alpha=0$, and the repeated-maximum conclusion of Theorem
\ref{thm:main} forces $\delta\le\gamma$. It is again nonnegative.
For $p=3$ or $p\equiv2\pmod3$, the valuation pattern is either
$(0,0,0,0)$ or a permutation of $(0,E,E,E)$, and this exponent is
zero in both cases. If the exponent is positive for a split prime, then
$\gamma>\delta$; the maximum must therefore occur at least twice
among the first three valuations. Thus $\beta=\gamma>0$.
\end{proof}

Rearranging the formula for $x$ gives precisely the algebraic part of Hohn's
fourth-radius comment. It makes no assertion about frequencies of the
possible values of $x$ or their prime factors.

\begin{corollary}\label{cor:bounds}
For a solution with $r_1\le r_2\le r_3\le r_4$, one has
$r_2\le3r_1$ and $r_3\le4r_1$. Given $r_1,r_2,r_3$, there is at most
one possible $r_4\ge r_3$. Consequently, only finitely many rows have a
specified first radius.
\end{corollary}
\begin{proof}
Set $X=1/r_1$, $Y=1/r_2$, $Z=1/r_3$, $W=1/r_4$, so
$X\ge Y\ge Z\ge W>0$, and put $T=Y+Z+W$. Cauchy's inequality and
\eqref{eq:radii} give
$(X+T)^2\ge3X^2+T^2$, hence $X\le T\le3Y$.
Let $F(X,Y,Z,W)=(X+Y+Z+W)^2-3(X^2+Y^2+Z^2+W^2)$.
For $0<W\le Z$ its derivative in $W$ is positive. Therefore
\[
 0=F(X,Y,Z,W)\le F(X,Y,Z,Z)
 =6XZ-\tfrac32X^2-2(Y-Z-X/2)^2,
\]
which implies $Z\ge X/4$. Finally, the two roots of the quadratic
equation for $W$ have sum $X+Y+Z>2Z$, so they cannot both lie in
$(0,Z]$. These statements give the claimed radius bounds, uniqueness,
and finiteness.
\end{proof}

\begin{corollary}\label{cor:membership}
Membership of a positive integer $N$ in A389940 is decidable by a finite
exact calculation. The calculation can also list every primitive sphere
quadruple containing $N$.
\end{corollary}
\begin{proof}
Any such ordered quadruple has $r_1\le N$, so Corollary
\ref{cor:bounds} bounds $r_2\le3N$ and $r_3\le4N$.
For each triple, put $a=r_2r_3$, $b=r_1r_3$, $c=r_1r_2$,
$s=a+b+c$, and $D=3(s^2-2(a^2+b^2+c^2))$.
The possible fourth radii are $2r_1r_2r_3/(s\pm\sqrt D)$.
Keep positive integral values at least $r_3$ with common gcd one and a
coordinate equal to $N$. A rational fourth radius makes $\sqrt D$
rational, hence integral; thus integer square tests suffice.
The finite list is empty exactly when $N$ is absent. This is an exhaustive
membership procedure, not a direct formula for terms. A389940, COMMENTS,
paragraph 3, asks whether terms can ``otherwise'' be found algorithmically;
a non-enumerative interpretation of that question is not answered here.
\end{proof}

\begin{corollary}\label{cor:chain}
There is no infinite nondecreasing sequence of positive rational radii
whose consecutive quadruples are distinct sphere solutions of
\eqref{eq:radii}. For real radii, any such sequence has constant ratio
$\lambda$, where $\lambda>1$ is the largest real root of
$z^4-z^3-z^2-z+1$ (the constant in A289917).
\end{corollary}
\begin{proof}
Write $b_n=1/r_n$. For a shared triple the two possible curvatures
sum to that triple's curvature sum, so distinct consecutive windows give
$b_{n+4}=b_{n+1}+b_{n+2}+b_{n+3}-b_n$.
The nonincreasing positive $b_n$ have a limit, and the recurrence makes
that limit zero. Its characteristic polynomial has four distinct roots:
$\lambda,\lambda^{-1}$ and two nonreal roots of modulus one. Indeed,
division by $z^2$ reduces it to $s^2-s-3=0$, $s=z+z^{-1}$;
one root $s$ exceeds $2$ and the other lies in $(-2,2)$.
The standard general solution of this constant-coefficient recurrence
therefore has only the $\lambda^{-n}$ component if it tends to zero:
the growing component is excluded by boundedness, and a nonzero
combination of the two unit-modulus components cannot tend to zero
(its mean squared modulus is positive).
Thus $b_n=C\lambda^{-n}$ with $C>0$.
The rational-root test excludes rational $\lambda$, so rational
radii are impossible. This proves the infinite-chain obstruction
asserted in A390148, COMMENTS, paragraph 3, but not its conjecture
of arbitrarily long finite primitive integer chains.
\end{proof}

\section{Repeated sphere radii}
Write $\omega(x)$ for the number of distinct prime factors of $x$, with
$\omega(1)=0$, and $\epsilon_3(x)=1$ if $3\mid x$, and $0$ otherwise.
These are A001221 and A079978 on positive arguments.
\begin{corollary}\label{cor:repeat}
Every repeated radius is $3x^2$ for some positive integer $x$.
For each $x$ there are exactly
\[
 2^{\omega(x)-\epsilon_3(x)+1}
\]
ordered primitive sphere quadruples containing that radius at least twice.
Explicitly, take $uv=x$, $\gcd(u,v)=1$, $3\nmid u$, and either sign in
\[
 A=u^2+3v^2\pm3uv.
\]
Sort the four numbers
\begin{equation}\label{eq:repeat}
 (3u^2v^2,\ 3u^2v^2,\ 3Av^2,\ Au^2).
\end{equation}
All the resulting quadruples are distinct. The only triple repetition
is $(1,3,3,3)$.
\end{corollary}
\begin{proof}
Repeated radii give repeated primitive curvatures, say $(A,A,B,C)$.
Equation \eqref{eq:integral} becomes
$(A-B-C)^2=3BC$. Interchange $B,C$ so that $v_3(B)$ is even and
$v_3(C)$ is odd. Prime factorization gives
$B=d u^2$, $C=3d v^2$ for a squarefree $d$ coprime to $3$.
Then $A=d(u^2+3v^2\pm3uv)$. Primitivity forces $d=1$,
$\gcd(u,v)=1$, and $3\nmid u$; otherwise a common prime divides
$A,B,C$. Conversely these conditions make $A,B,C$ pairwise coprime:
modulo a prime dividing $u$, $A\equiv3v^2$, and modulo a prime
dividing $3v$, $A\equiv u^2$.
Their least common multiple is $3Au^2v^2$, giving \eqref{eq:repeat}.
The minus sign also gives $A>0$, since
$u^2-3uv+3v^2=(u-3v/2)^2+3v^2/4$.

Each prime-power factor of $x$ must belong wholly to $u$ or to $v$;
the factor at $3$, if present, must belong to $v$. There are
$2^{\omega(x)-\epsilon_3(x)}$ choices and two distinct signs.
There is no overcounting: primitive radii uniquely recover primitive
curvatures, their repeated coordinate $A$, and the remaining pair
$\{B,C\}$; its square and three-times-square members distinguish $u,v$.
If a third curvature equals $A$, either $A=u^2$ or $A=3v^2$.
The first forces the minus sign and $u=v=1$; the second would force
$u=3v$, contrary to $3\nmid u$. Thus even the triple-repetition case
has a unique parameter choice. Four equal curvatures do not satisfy
\eqref{eq:integral}.
\end{proof}

\section{The linked circle conjecture}
Hohn's linked array A390583 \cite{circles} uses coefficient $2$ in
\eqref{eq:radii}, not $3$. Its COMMENTS, paragraph 3, conjectures
a cube-or-twice-a-cube difference when the largest two radii coincide.
The following classification proves it, even without the largest-radius
restriction.
\begin{corollary}\label{cor:circle}
Every primitive positive circle quadruple with repeated radii and
distinct remaining radii has the form, up to permutation,
\[
 (a,b,c,c)=\frac1h
 (u d^2,\ v d^2,\ 4uv(u+v),\ 4uv(u+v)),
 \quad d=v-u>0,\quad \gcd(u,v)=1,
\]
where $h=\gcd(d^2,4(u+v))$. Consequently $b-a$ is a positive cube
or twice a positive cube.
\end{corollary}
\begin{proof}
Put $X=1/a$, $Y=1/b$, $Z=1/c$. The circle equation gives
$(X-Y)^2=4Z(X+Y)$, so $a=b$ is impossible and
$c=4ab(a+b)/(b-a)^2$. Writing $a/b=u/v$ in lowest terms yields
the displayed vector after clearing denominators and dividing by its
common gcd. Since $\gcd(d,uv)=1$, that gcd is precisely $h$.
No odd prime divides $h$, because an odd prime dividing both $d$ and
$u+v=2u+d$ would divide $u$ and $v$. If $d$ is odd, $h=1$.
If $v_2(d)=1$, $u,v$ are odd, $4\mid u+v$, and $h=4$.
If $v_2(d)\ge2$, $v_2(u+v)=1$ and $h=8$.
Thus $b-a=d^3/h$ is respectively $d^3$, $2(d/2)^3$, or
$(d/2)^3$. The formula also proves the converse parametrization.
\end{proof}

\begin{corollary}\label{cor:circlebounds}
Ordered positive circle quadruples satisfy $r_1<r_2$,
$r_2\le(3+2\sqrt3)r_1$, and
\[
 r_3\le\frac{4(1/r_1+1/r_2)}{(1/r_1-1/r_2)^2}.
\]
There is at most one $r_4\ge r_3$ for each triple, hence finitely
many rows for each first radius. The following three repeated-radius
families hold, with the second correcting the index in
A390583, COMMENTS, paragraph 4:
\[
\begin{split}
 &(k,k+1,4k(k+1)(2k+1),4k(k+1)(2k+1)),\\
 &(k,k+2,2k(k+1)(k+2),2k(k+1)(k+2))\\
 &\hspace{35mm}\text{for odd }k,\\
 &(k,k+8,k(k+4)(k+8)/8,k(k+4)(k+8)/8)\\
 &\hspace{35mm}\text{for }k\equiv2\pmod4.
\end{split}
\]
\end{corollary}
\begin{proof}
Write $X\ge Y\ge Z\ge W>0$ for reciprocal radii and
$G=(X+Y+Z+W)^2-2(X^2+Y^2+Z^2+W^2)$.
If $X=Y$, then $G=4X(Z+W)-(Z-W)^2>0$, impossible.
The derivative of $G$ in each of $Y,Z,W$, when that coordinate is
at most $Y$, is positive. Hence
$0=G\le G(X,Y,Y,Y)=-X^2+6XY+3Y^2$,
giving the first bound. Also
$0\le G(X,Y,Z,Z)=4Z(X+Y)-(X-Y)^2$, giving the second.
The two roots for $W$ have sum $2(X+Y+Z)>2Z$, so at most one
lies in $(0,Z]$. Substitute $b-a=1,2,8$ in
$c=4ab(a+b)/(b-a)^2$ for the three families. They are integral
and primitive: the first two have $\gcd(a,b)=1$, while in the
third $\gcd(a,b)=2$ and $c$ is odd. Each has $c\ge b$.
\end{proof}

The source writes the second family's repeated radius as
``A000447(k)*12''. The defining formula is
\[
 \mathrm{A000447}(n)=n(4n^2-1)/3.
\]
At $k=3$ the stated index gives $(3,5,420,420)$, which fails the
circle equation:
\[
 \left(\sum 1/r_i\right)^2-2\sum 1/r_i^2=-4/315.
\]
The correct index is $(k+1)/2$. It gives repeated radius
$2k(k+1)(k+2)$, hence $(3,5,120,120)$ at $k=3$.
The first and third families agree with the stated A069074 and A061550
indices, respectively \cite{circlefunctions}.

\section{Coverage and verification}
The following table distinguishes the mathematical assertions in the posing
entry from its numerical observations and other questions.
\begin{longtable}{@{}>{\raggedright\arraybackslash}p{0.27\linewidth}>{\raggedright\arraybackslash}p{0.20\linewidth}>{\raggedright\arraybackslash}p{0.46\linewidth}@{}}
\toprule
Source item & Coverage & Statement\\
\midrule
\endhead
A390148, C5 & Prior proof cited & Kontorovich's classification implies the $3$-adic pattern; Theorem \ref{thm:main} supplies a direct proof.\\[3pt]
A390148, C6--C7 & Proved; prior arithmetic cited & Theorem \ref{thm:main}: all-prime thresholds and inert-prime pattern; Stong's norm argument gives the latter.\\[3pt]
A390148, C9 & Proved in part & Corollary \ref{cor:gcd}: integrality and split-prime support. Frequency clauses: not covered.\\[3pt]
A390148, C2 & Proved & Corollary \ref{cor:bounds}: both bounds, uniqueness, and finite rows per first radius.\\[3pt]
A390148, C8 & Proved & Corollary \ref{cor:repeat}: exact count and classification of every repeated radius.\\[3pt]
A390148, C3 & Proved in part & Corollary \ref{cor:chain}: no infinite rational chain. Arbitrarily long finite chains: not covered.\\[3pt]
A390322, A388858, A390464, A390465 & Proved consequences & The four column definitions transfer the rowwise results; no column-only characterization is asserted.\\[3pt]
A389940, C3: algorithm & Not covered as a direct construction & Corollary \ref{cor:membership} gives exhaustive finite membership (not posed; proved).\\[3pt]
A389940, C2: examples & Finite check & The only containing rows for $1,4,36$ are the three rows listed there.\\[3pt]
A389940, C3: gaps & Not covered & No asymptotic gap-growth statement.\\[3pt]
A001221, A079978 & Definitions used & $\omega(x)$ and $\epsilon_3(x)$ in Corollary \ref{cor:repeat}.\\[3pt]
A289917 & Proved consequence & Corollary \ref{cor:chain}: its quartic root is the only infinite real-chain ratio. Other formulas: not covered.\\[3pt]
A390583, C2--C3 & Proved & Corollaries \ref{cor:circle}--\ref{cor:circlebounds}: cubic difference, bounds and uniqueness.\\[3pt]
A390583, C4 & Refuted-and-corrected in part & Second-family index corrected by explicit example and formula; first and third families proved.\\[3pt]
A069074, A000447, A061550 & Formulas used & Corollary \ref{cor:circlebounds} identifies the three polynomial families, with the A000447 index corrected.\\[3pt]
A290508 & Not posed; proved & Corollary \ref{cor:circle} classifies its repeated-largest-radius subfamily; the general largest-radius sequence is not characterized.\\[3pt]
A390583, C5 & Not covered & The arbitrarily-long-chain question remains outside these results.\\
\bottomrule
\end{longtable}
Here C$n$ means the $n$th COMMENTS paragraph in the cited version.
A390148 uses revision 23; A389940 uses revision 13 (5 November 2025);
A390583 uses revision 10 (17 November 2025). Statements described as
not covered are not counterexamples to their sources.

\paragraph{Lean statements.}
The declaration
\href{https://github.com/the-omega-institute/trureturing/blob/bed437cb2aa5a2d3570eb2241767ad6de36d77cc/D5/S3/Arith/Descartes/SplitPrimeSphereRadii.lean}{\texttt{SplitPrimeSphereRadii.result}}
has exactly the positive, common-primitive, coefficient-three hypotheses
above and the threshold-count conclusion for $p\equiv1\pmod3$.
The declaration
\href{https://github.com/the-omega-institute/trureturing/blob/bed437cb2aa5a2d3570eb2241767ad6de36d77cc/D5/S3/Arith/Descartes/PrimitiveSphereRadii.lean}{\texttt{PrimitiveSphereRadii.primitive\_sphere\_radii\_v3}}
has the $3$-adic pattern as its conclusion. The all-prime extension,
inert-prime parity, and all corollaries are established by the written
proofs above. Neither declaration asserts those additional conclusions.
Both links pin revision \texttt{bed437cb2aa5a2d3570eb2241767ad6de36d77cc}.

\paragraph{Exact finite checks.}
The accompanying \texttt{verify.py} checks the reciprocal equation,
primitivity, valuation patterns, all positive divisibility thresholds, and
the gcd correction on all $1128$ quadruples in Hohn's $4512$-term table.
It also enumerates the exact quadratic solutions in the bounds of
Corollary \ref{cor:bounds} for $1\le r_1\le60$, obtaining the same $19$
rows as the table. The containing-row checks for $1,4,36$ are exhaustive
by Corollary \ref{cor:membership}. Further checks test the repeated-sphere
parametrization for $x\le200$ and the circle parametrization for coprime
$1\le u<v\le100$, including the three displayed circle families.
All operations use integers or rational numbers.
The input digest, ranges, and detailed counts are in
\texttt{verification.json}; these finite checks illustrate the unbounded
written proofs, rather than replacing them.

\begin{thebibliography}{10}
\bibitem{hohn}
Charles L. Hohn, \emph{Entry A390148}, The On-Line Encyclopedia of Integer Sequences,
revision 23, 17 November 2025, COMMENTS, the successive prime-divisibility
and gcd-formula assertions.
\url{https://oeis.org/A390148}.
The exact revision is accessible from
\url{https://oeis.org/history/view?seq=A390148&v=23}.

\bibitem{circles}
Charles L. Hohn, \emph{Entry A390583}, The On-Line Encyclopedia of Integer
Sequences, revision 10, 17 November 2025, COMMENTS, paragraphs 2--5.
\url{https://oeis.org/A390583}.

\bibitem{values}
Charles L. Hohn, \emph{Entry A389940}, The On-Line Encyclopedia of Integer
Sequences, revision 13, 5 November 2025, COMMENTS, paragraphs 2--3.
\url{https://oeis.org/A389940}.

\bibitem{auxiliary}
The On-Line Encyclopedia of Integer Sequences, entries
\href{https://oeis.org/A001221}{A001221} and
\href{https://oeis.org/A079978}{A079978} (arithmetic functions),
\href{https://oeis.org/A289917}{A289917} (spiral ratio), and
\href{https://oeis.org/A390322}{A390322},
\href{https://oeis.org/A388858}{A388858},
\href{https://oeis.org/A390464}{A390464},
\href{https://oeis.org/A390465}{A390465} (columns).

\bibitem{lmw}
Jeffrey C. Lagarias, Colin L. Mallows, and Allan R. Wilks,
\emph{Beyond the Descartes circle theorem},
American Mathematical Monthly \textbf{109} (2002), no.~4, 338--361.
\href{https://doi.org/10.1080/00029890.2002.11920896}{doi:10.1080/00029890.2002.11920896}.
Theorem 1.2 in the author version,
\url{https://arxiv.org/abs/math/0101066v1}.

\bibitem{kontorovich}
Alex Kontorovich, \emph{The local-global principle for integral Soddy
sphere packings}, arXiv:1208.5441v1, 27 August 2012, Lemma 2.11,
p.~9. \url{https://arxiv.org/abs/1208.5441v1}.
Journal version: Journal of Modern Dynamics \textbf{15} (2019), 209--236,
\href{https://doi.org/10.3934/jmd.2019019}{doi:10.3934/jmd.2019019}.

\bibitem{mcgonagle}
Annmarie McGonagle and Sam Northshield,
\emph{A new parameterization of Ford circles}, Pi Mu Epsilon Journal
\textbf{13} (2014), no.~10, 637--643, Section 6.
\href{https://www.researchgate.net/publication/329218164_A_New_Parameterization_of_Ford_Circles}{Author manuscript},
dated 12 November 2013, Section 6, pp.~8--9.

\bibitem{stong}
Richard Stong, solution of Sam Northshield's Problem 11671,
\emph{If the sum of the squares is the square of the sum, \ldots},
American Mathematical Monthly \textbf{122} (2015), no.~1, p.~76
(additional solvers, p.~77).
\href{https://doi.org/10.4169/amer.math.monthly.122.01.75}{doi:10.4169/amer.math.monthly.122.01.75}.

\bibitem{northshield}
Sam Northshield, \emph{Ford circles and spheres},
arXiv:1503.00813v1, 3 March 2015, Theorem 4.4 and Corollary 4.6,
pp.~13--14. \url{https://arxiv.org/abs/1503.00813v1}.

\bibitem{circlefunctions}
The On-Line Encyclopedia of Integer Sequences, entries
\href{https://oeis.org/A069074}{A069074},
\href{https://oeis.org/A000447}{A000447},
\href{https://oeis.org/A061550}{A061550}, and
\href{https://oeis.org/A290508}{A290508}.
\end{thebibliography}

\medskip
\noindent\small Copyright 2026 Haobo Ma and Wenlin Zhang.
Licensed under the Creative Commons Attribution 4.0 International License
(\href{https://creativecommons.org/licenses/by/4.0/}{CC BY 4.0}).
\end{document}
